
%%=============================================================================
%% ESTADO DEL ARTE
%%=============================================================================
\chapter{Estado del arte}

La seguridad de los sistemas embebidos y del Internet de las Cosas constituye actualmente una de las principales áreas de investigación dentro de la ciberseguridad moderna. El crecimiento exponencial del número de dispositivos conectados \cite{Sinha2025StateGlobally}, junto con el avance de la computación distribuida, el \textit{edge computing} y las ciudades inteligentes, ha incrementado significativamente la necesidad de desarrollar mecanismos criptográficos capaces de operar bajo severas restricciones de memoria, capacidad de procesamiento y consumo energético.

Durante los últimos años, la comunidad científica ha concentrado sus esfuerzos en tres líneas principales de investigación. La primera corresponde al diseño de algoritmos de criptografía ligera que permitan proteger la información utilizando una cantidad reducida de recursos computacionales. La segunda se enfoca en el desarrollo de mecanismos resistentes frente a la futura amenaza de la computación cuántica. Finalmente, una tercera línea de investigación, menos explorada pero igualmente crítica, estudia la generación de números aleatorios de alta calidad como fundamento de cualquier sistema criptográfico seguro.

El presente capítulo revisa el estado del conocimiento en estas tres áreas, iniciando con una descripción de la evolución del Internet de las Cosas y sus restricciones computacionales, continuando con los principales desafíos de seguridad presentes en estos entornos. Posteriormente, se analiza la evolución de la criptografía ligera y de los mecanismos modernos de generación de entropía, para finalmente estudiar los fundamentos del \textit{Quantum Permutation Pad} (QPP) y \textit{Quantum Permutation Pad Random Number Generator} (QPP-RNG), identificando las brechas existentes en la literatura que motivan la presente investigación.


%%=========== 
%% SECCION 1
%%=========== 

\section{Internet de las cosas: evolución, arquitectura y restricciones}

El concepto de Internet de las Cosas (\textit{Internet of Things}, IoT) hace referencia a la interconexión de objetos físicos capaces de percibir información del entorno mediante sensores, procesarla localmente y transmitirla a través de diferentes tecnologías de comunicación hacia otros dispositivos o plataformas de procesamiento. Aunque el término fue propuesto por Kevin Ashton a finales de la década de los noventa \cite{Dhanda2020LightweightIoT}, solamente durante la última década, la disponibilidad de sensores de bajo costo, redes inalámbricas, servicios en la nube y plataformas de análisis masivo de datos permitieron su adopción a gran escala.

Actualmente el IoT constituye uno de los principales habilitadores tecnológicos de la Industria 4.0. Este permite integrar procesos físicos y digitales mediante arquitecturas distribuidas donde  millones de dispositivos trabajan de la mano para monitorear y controlar procesos en tiempo real. Esta evolución ha favorecido el desarrollo de aplicaciones en múltiples dominios, incluyendo ciudades inteligentes, automatización industrial, agricultura de precisión y sistemas de transporte inteligente, entre otros.
Desde una perspectiva arquitectónica, la mayoría de implementaciones IoT pueden describirse mediante un modelo jerárquico compuesto por cuatro niveles principales.

El primer nivel corresponde a la capa de percepción, integrada por sensores, actuadores y dispositivos embebidos encargados de capturar información del entorno físico. Estos dispositivos representan el punto de conexión entre el mundo físico y el mundo digital. Generalmente poseen capacidades computacionales extremadamente limitadas.

El segundo nivel corresponde a la capa de comunicación, responsable del transporte seguro de la información utilizando tecnologías como Wi-Fi, Bluetooth Low Energy (BLE), Zigbee, Thread, LoRaWAN, NB-IoT, Ethernet industrial y redes celulares de quinta generación (5G). Cada tecnología presenta ventajas particulares en términos de alcance, ancho de banda, latencia y consumo energético, condicionando directamente las soluciones criptográficas que pueden implementarse.

En tercer lugar, la capa de procesamiento incorpora nodos Edge, Fog o servicios Cloud donde se ejecutan procesos de almacenamiento, agregación de datos, análisis estadístico, inteligencia artificial y automatización.

Finalmente, la capa de aplicación implementa la lógica de negocio específica para cada dominio, proporcionando servicios a usuarios finales, operadores industriales o sistemas de apoyo a la toma de decisiones.

Uno de los aspectos que diferencia al Internet de las Cosas de los sistemas informáticos tradicionales radica en las severas restricciones de recursos presentes en la mayoría de dispositivos embebidos. Mientras un computador personal moderno dispone de varios gigabytes de memoria RAM, múltiples núcleos de procesamiento y almacenamiento de estado sólido, gran parte de los microcontroladores utilizados en aplicaciones IoT operan con frecuencias inferiores a los 100 MHz, memorias RAM del orden de pocos kilobytes y capacidades de almacenamiento inferiores a un megabyte.

Estas limitaciones afectan directamente la implementación de mecanismos criptográficos convencionales. Algoritmos ampliamente utilizados como RSA, ECC o incluso implementaciones completas de AES pueden consumir una fracción considerable de la memoria disponible o incrementar significativamente el tiempo de procesamiento requerido para cada operación criptográfica. De igual manera, muchos dispositivos funcionan mediante baterías cuya autonomía debe extenderse durante varios años, haciendo indispensable minimizar el número de operaciones matemáticas ejecutadas y reducir al máximo el consumo energético asociado a los procesos de autenticación y cifrado. Estudios recientes muestran que incluso pequeños incrementos en la complejidad computacional pueden traducirse en reducciones significativas de la vida útil de los dispositivos desplegados en campo \cite{Radhakrishnan2024EfficiencyDevices}.

Las restricciones descritas anteriormente han motivado la aparición del concepto de dispositivos con recursos restringidos (\textit{resource-constrained devices}). Diversos organismos internacionales, incluyendo el Internet Engineering Task Force (IETF) y el National Institute of Standards and Technology (NIST), clasifican estos dispositivos según su capacidad de memoria, procesamiento y energía disponible \cite{McKay2017ReportCryptography}. En particular, los dispositivos pertenecientes a las clases más bajas pueden disponer únicamente de algunos kilobytes de memoria RAM y capacidades de procesamiento limitadas, haciendo prácticamente imposible la ejecución de protocolos criptográficos diseñados originalmente para computadores.

Como consecuencia, el diseño de soluciones de seguridad para IoT requiere optimizar simultáneamente múltiples variables: tiempo de ejecución, utilización de memoria, consumo energético, tamaño del código ejecutable, ancho de banda requerido para transmisión y resistencia frente a ataques criptográficos. Esta combinación de restricciones convierte al Internet de las Cosas en uno de los escenarios más exigentes para el diseño de algoritmos criptográficos modernos.

Otro aspecto diferenciador del IoT consiste en la heterogeneidad de sus componentes y características físicas. En una misma arquitectura pueden coexistir microcontroladores ARM Cortex-M, procesadores RISC-V, ESP32, AVR, FPGA e incluso computadores industriales de altas capacidades. Esta diversidad dificulta el desarrollo de bibliotecas criptográficas portables y obliga a diseñar implementaciones suficientemente flexibles para adaptarse a múltiples arquitecturas sin comprometer la seguridad ni el rendimiento.

Adicionalmente, la naturaleza distribuida de estos sistemas incrementa la complejidad de la gestión de claves criptográficas, la autenticación mutua entre dispositivos y la actualización segura del software. Miles de dispositivos pueden permanecer desplegados durante años en ubicaciones remotas donde las tareas de mantenimiento físico resultan costosas o incluso inviables. Bajo estas condiciones, cualquier vulnerabilidad criptográfica puede permanecer explotable durante largos periodos de tiempo, afectando simultáneamente un gran número de nodos dentro de la infraestructura.

En consecuencia, la comunidad científica coincide en que la seguridad del Internet de las Cosas no puede abordarse simplemente trasladando las soluciones desarrolladas para sistemas tradicionales. Las restricciones computacionales, la heterogeneidad del hardware, la escalabilidad de las implementaciones y la necesidad de minimizar el consumo energético requieren mecanismos específicamente diseñados para estos escenarios. Esta necesidad dio origen a una nueva línea de investigación conocida como Criptografía Ligera (Lightweight Cryptography), cuyo objetivo consiste en proporcionar niveles elevados de seguridad manteniendo simultáneamente un consumo reducido de recursos computacionales \cite{McKay2017ReportCryptography,NationalInstituteofStandards2023AnnouncingSelection}. Dicho paradigma constituye actualmente una de las principales áreas de investigación dentro de la seguridad informática aplicada a sistemas embebidos e infraestructura crítica.


%%=========== 
%% SECCION 2
%%=========== 

\section{Seguridad en dispositivos IoT}

La seguridad constituye uno de los principales retos para la adopción masiva del Internet de las Cosas. A diferencia de los sistemas informáticos convencionales, los dispositivos IoT operan en entornos altamente distribuidos, heterogéneos y en muchos casos, físicamente accesibles para atacantes potenciales. Esta combinación de restricciones computacionales, conectividad permanente y despliegues de gran escala ha incrementado significativamente la superficie de ataque de las infraestructuras digitales modernas, convirtiendo la protección de estos sistemas en una prioridad tanto para la comunidad científica como para organismos gubernamentales y el sector industrial.

La mayoría de los dispositivos IoT son diseñados para desempeñar funciones específicas durante largos periodos de tiempo con una mínima intervención humana. En consecuencia, suelen permanecer conectados de manera continua a redes públicas o privadas, transmitiendo información sensible relacionada con procesos industriales, variables ambientales, datos médicos o infraestructura crítica. Cualquier vulnerabilidad presente en estos dispositivos puede comprometer no solamente la confidencialidad de la información, sino también la integridad de los datos, la disponibilidad de los servicios e incluso la seguridad física de personas e instalaciones.

Diversos estudios coinciden en que el modelo de amenazas presente en IoT difiere considerablemente del observado en arquitecturas tradicionales \cite{Seyhan2022ClassificationTaxonomy}. Mientras los sistemas corporativos suelen estar protegidos mediante múltiples capas de seguridad perimetral, los dispositivos IoT frecuentemente operan en ubicaciones remotas, redes inalámbricas y ambientes físicamente inseguros, donde el atacante puede tener acceso tanto al canal de comunicación como al propio dispositivo. Estas condiciones obligan a considerar escenarios de ataque significativamente más amplios, incluyendo ataques físicos, ataques de software, ataques criptográficos y ataques sobre los mecanismos de generación de números aleatorios.

\subsection{Modelo de amenazas en entornos IoT}

El diseño de mecanismos de seguridad para dispositivos IoT parte del establecimiento de un modelo de amenazas que permita identificar los activos críticos, los posibles atacantes y las capacidades que estos podrían poseer. Generalmente se asume que un atacante puede interceptar, modificar o reproducir el tráfico de red, comprometer dispositivos individuales, acceder físicamente a sensores desplegados en campo e incluso ejecutar ataques coordinados sobre múltiples nodos pertenecientes a una misma infraestructura.

Dentro de este modelo, los activos que requieren protección incluyen las claves criptográficas, credenciales de autenticación, firmware, datos almacenados localmente, información transmitida a través de la red y mecanismos de actualización remota del software. La pérdida de cualquiera de estos activos puede comprometer completamente la seguridad del sistema.

Desde la perspectiva del atacante, los objetivos más comunes consisten en obtener acceso no autorizado a la información, modificar el comportamiento del dispositivo, suplantar su identidad, alterar mediciones provenientes de sensores o utilizar el dispositivo como punto de entrada hacia otros componentes de la infraestructura. Este último escenario ha adquirido especial relevancia debido al incremento de ataques distribuidos donde miles de dispositivos IoT comprometidos son utilizados para ejecutar campañas de denegación de servicio (Distributed Denial of Service, DDoS) \cite{Klimushyn2022Crypto-ResistantDevices}.

\subsection{Ataques físicos}

Una característica particular del Internet de las Cosas consiste en que muchos dispositivos permanecen desplegados en ambientes físicamente accesibles para potenciales atacantes. Sensores ambientales, medidores inteligentes, dispositivos agrícolas y equipos industriales suelen instalarse en lugares donde resulta imposible garantizar protección física permanente.

Esta situación permite la ejecución de ataques que difícilmente serían viables sobre servidores tradicionales. Entre ellos se encuentran la extracción directa de memoria Flash, lectura de buses internos mediante interfaces JTAG o SWD, manipulación de voltajes de alimentación, ataques por inyección de fallos (\textit{fault injection}), análisis mediante canales laterales (\textit{side-channel attacks}) y técnicas de ingeniería inversa sobre el firmware \cite{Klimushyn2022Crypto-ResistantDevices}.

Los ataques de canal lateral representan una amenaza especialmente importante debido a que permiten recuperar información criptográfica analizando variables indirectas como el consumo de energía, emisiones electromagnéticas o tiempos de ejecución. Incluso cuando el algoritmo criptográfico utilizado sea matemáticamente seguro, una implementación deficiente puede filtrar suficiente información para permitir la reconstrucción de claves secretas.

\subsection{Ataques sobre los canales de comunicación}

La comunicación inalámbrica constituye otro de los principales puntos de vulnerabilidad dentro de las arquitecturas IoT. Tecnologías como Wi-Fi, Bluetooth Low Energy, Zigbee, Thread o LoRaWAN utilizan canales compartidos susceptibles a interceptación, interferencia y manipulación por parte de terceros.

Entre los ataques más frecuentes se encuentran la escucha pasiva (\textit{eavesdropping}), los ataques de repetición (\textit{replay attacks}), la modificación de mensajes en tránsito, la suplantación de dispositivos (\textit{spoofing}), los ataques de hombre en la mitad (\textit{Man-in-the-Middle}) y las campañas de denegación de servicio.

La naturaleza distribuida del IoT incrementa además la dificultad para establecer mecanismos robustos de autenticación mutua entre dispositivos. Un atacante capaz de introducir un nodo falso dentro de la red puede manipular el flujo de información, alterar mediciones provenientes de sensores o inyectar comandos maliciosos dirigidos hacia actuadores industriales.

\subsection{Ataques criptográficos}

Los mecanismos criptográficos constituyen la principal línea de defensa frente a los ataques descritos anteriormente. Sin embargo, la seguridad de un sistema criptográfico depende simultáneamente de la robustez matemática del algoritmo, de la correcta implementación del software y de la calidad de las claves utilizadas durante el proceso de cifrado.

Entre los ataques criptográficos más relevantes para dispositivos IoT se encuentran el análisis diferencial, el análisis lineal, los ataques por texto conocido, ataques por texto elegido, ataques sobre protocolos de autenticación y ataques dirigidos contra la gestión de claves.

La llegada de la computación cuántica introduce además nuevas amenazas para los algoritmos asimétricos tradicionales. Algoritmos como RSA y la criptografía basada en curvas elípticas dependen de problemas matemáticos que podrían resolverse eficientemente mediante el algoritmo de Shor una vez existan computadores cuánticos suficientemente grandes \cite{Shor1994AlgorithmsFactoring}. Los algoritmos simétricos, a pesar de no ser vulnerables ante el algoritmo de Shor, ven reducida a la mitad su fortaleza efectiva frente al algoritmo de búsqueda de Grover \cite{Grover1996ASearch}. Aunque los algoritmos simétricos presentan una mayor resistencia frente a este escenario, continúan requiriendo claves de alta calidad y fuentes confiables de entropía para mantener sus propiedades de seguridad \cite{Turan2018RecommendationGeneration}.

\subsection{Gestión de claves criptográficas}

La administración de claves representa uno de los mayores desafíos dentro de las infraestructuras IoT. Los dispositivos deben intercambiar información de manera segura durante largos periodos de tiempo, frecuentemente utilizando redes poco confiables y con capacidad limitada para realizar procesos complejos de autenticación.

Una gestión inadecuada de claves puede comprometer completamente la seguridad del sistema, independientemente de la fortaleza del algoritmo criptográfico utilizado. Entre los problemas más frecuentes se encuentran la reutilización de claves, almacenamiento inseguro, distribución mediante canales no protegidos, ausencia de mecanismos de rotación periódica y utilización de credenciales predeterminadas.

En consecuencia, numerosos trabajos recientes enfatizan que la seguridad de un sistema criptográfico depende tanto del algoritmo de cifrado como de la calidad del material criptográfico empleado durante la generación de claves, nonces, semillas y vectores de inicialización \cite{MITRE2024CWE-327:Algorithm}.

\subsection{La importancia de las fuentes de entropía}

Uno de los componentes menos visibles pero más importantes dentro de cualquier sistema criptográfico corresponde al mecanismo encargado de generar números aleatorios. Claves criptográficas, nonces, semillas para generadores pseudoaleatorios, protocolos de autenticación y numerosos mecanismos de seguridad dependen directamente de la disponibilidad de secuencias impredecibles.

La historia de la criptografía moderna registra múltiples incidentes donde algoritmos matemáticamente seguros fueron comprometidos debido a fuentes deficientes de entropía \cite{Menezes1996HandbookCryptography,Klimushyn2022Crypto-ResistantDevices}. Generadores de números aleatorios con baja variabilidad pueden producir claves repetidas o parcialmente predecibles, facilitando ataques que reducen significativamente el espacio de búsqueda necesario para recuperar información confidencial.

Este problema resulta especialmente crítico en dispositivos IoT, donde las restricciones de hardware limitan la disponibilidad de fenómenos físicos capaces de producir verdadera aleatoriedad. Muchos microcontroladores carecen de generadores físicos de números aleatorios, obligando a obtener entropía a partir del comportamiento interno del sistema operativo o de la microarquitectura del procesador. Como consecuencia, durante los últimos años ha surgido un creciente interés por desarrollar mecanismos capaces de producir niveles adecuados de entropía sin requerir hardware especializado, entre ellos Quantum Permutation Pad Random Number Generator (QPP-RNG) \cite{Kuang2025QPP-RNG:Randomness,Kuang2025IID-BasedJitter}, objeto central de la presente investigación.

\subsection{Discusión}

La revisión realizada evidencia que la seguridad del Internet de las Cosas constituye un problema multidimensional que trasciende la simple selección de un algoritmo criptográfico. Las restricciones computacionales, la exposición física de los dispositivos, la heterogeneidad del hardware y la creciente sofisticación de los ataques obligan a considerar soluciones de seguridad integrales.

En este contexto, la criptografía ligera emerge como una respuesta natural a las limitaciones de los dispositivos embebidos. Sin embargo, incluso los algoritmos ligeros más robustos dependen de la disponibilidad de fuentes de entropía capaces de generar material criptográfico de alta calidad. Esta observación motiva el estudio de nuevas propuestas como QPP-RNG, cuyo objetivo consiste en proporcionar mecanismos de generación de números aleatorios adecuados para plataformas con recursos restringidos, manteniendo simultáneamente elevados niveles de seguridad estadística y eficiencia computacional.

La siguiente sección analiza la evolución de la criptografía ligera, el proceso de estandarización liderado por el National Institute of Standards and Technology (NIST) y las principales propuestas desarrolladas para proteger dispositivos IoT con restricciones severas de memoria, procesamiento y consumo energético.

%%=========== 
%% SECCION 3
%%=========== 

\section{Criptografía ligera y estandarización para dispositivos restringidos}

La criptografía ligera, conocida en la literatura como \textit{Lightweight Cryptography} (LWC), surge como respuesta a la brecha existente entre los requisitos de seguridad de los sistemas modernos y las limitaciones físicas de los dispositivos embebidos. A diferencia de la criptografía tradicional, cuyo diseño suele estar orientado a computadores no especializados, servidores o infraestructuras con alta capacidad de procesamiento, la criptografía ligera busca proporcionar mecanismos de confidencialidad, integridad y autenticación utilizando una cantidad reducida de memoria, ciclos de CPU, ancho de banda y energía \cite{McKay2017ReportCryptography,Althobaiti2025LightweightStandardization,Thakor2021LightweightOpportunities}.

El objetivo de la criptografía ligera no consiste en reducir la seguridad, sino en optimizar la relación entre protección y costo computacional. Esta distinción resulta fundamental, pues un algoritmo ligero no debe entenderse como una versión ``débil'' de un algoritmo tradicional, sino como una solución diseñada específicamente para operar en contextos donde los recursos disponibles son escasos. En consecuencia, el diseño de primitivas criptográficas ligeras exige equilibrar simultáneamente la seguridad y resistencia frente a ataques conocidos, eficiencia en software, eficiencia en hardware, tamaño de código, latencia, consumo energético y facilidad de implementación.

En los entornos IoT, esta necesidad se vuelve especialmente relevante. Un dispositivo embebido puede disponer únicamente de algunos kilobytes de memoria RAM, una frecuencia de reloj reducida y una batería que debe operar durante meses o años sin reemplazo. Bajo estas condiciones, el uso de algoritmos convencionales puede generar sobrecostos significativos. Por ejemplo, esquemas asimétricos tradicionales como RSA requieren operaciones aritméticas de alta complejidad, mientras que implementaciones completas de algoritmos simétricos robustos pueden ocupar una proporción considerable de la memoria disponible. Por esta razón, la literatura reciente ha enfatizado la necesidad de diseñar algoritmos capaces de reducir el costo computacional sin comprometer las propiedades criptográficas esenciales.

\subsection{Familias de algoritmos de criptografía ligera}

La literatura clasifica los algoritmos de criptografía ligera en varias familias, dependiendo de su estructura interna y del tipo de servicio de seguridad que proporcionan. Entre las familias más relevantes se encuentran los cifradores de bloque, cifradores de flujo, algoritmos basados en redes SPN, construcciones ARX, funciones hash ligeras y esquemas de cifrado autenticado con datos asociados \cite{Dhanda2020LightweightIoT}.

Los cifradores de bloque ligeros procesan bloques de información de tamaño fijo mediante rondas sucesivas de transformación. Ejemplos representativos incluyen PRESENT, LED, SIMON y SPECK \cite{Bogdanov2007PRESENT:Cipher,Beaulieu2015TheCiphers}. Estos algoritmos suelen estar diseñados para reducir el número de compuertas en hardware o minimizar el tamaño de código en software, manteniendo un nivel de seguridad adecuado para aplicaciones de bajo consumo.

Las redes de sustitución-permutación, conocidas como SPN, constituyen una de las estructuras más utilizadas en criptografía simétrica. En estas construcciones, la confusión se obtiene mediante sustituciones no lineales y la difusión mediante permutaciones o mezclas lineales. AES es el ejemplo más reconocido de esta familia \cite{NationalInstituteofStandards2001AdvancedAES}, aunque no fue diseñado originalmente como algoritmo ligero. No obstante, debido a su amplio análisis criptográfico y madurez, AES-128 continúa utilizándose como punto de referencia para comparar nuevas propuestas.

Otra familia importante corresponde a los cifradores ARX, cuyo nombre proviene de las operaciones \textit{Addition}, \textit{Rotation} y \textit{XOR}. Estas operaciones son eficientes en procesadores de propósito general y microcontroladores, ya que suelen estar soportadas directamente por el conjunto de instrucciones del hardware. SPECK pertenece a esta familia y ha sido ampliamente estudiado debido a su simplicidad, bajo consumo de memoria y buen desempeño en plataformas restringidas \cite{Beaulieu2015TheCiphers}.

Finalmente, los esquemas AEAD, o \textit{Authenticated Encryption with Associated Data}, proporcionan simultáneamente confidencialidad, integridad y autenticación. Esta característica resulta especialmente importante en IoT, donde no basta con cifrar la información transmitida, sino que también es necesario garantizar que los mensajes no hayan sido modificados y que provengan de una fuente legítima. ASCON pertenece a esta categoría y actualmente constituye uno de los principales referentes de criptografía ligera para dispositivos restringidos \cite{NationalInstituteofStandards2023AnnouncingSelection,Turan2025Ascon-BasedFunctions}.

\subsection{Cifrado autenticado con datos asociados}

El cifrado autenticado con datos asociados representa una de las construcciones más importantes para la seguridad de sistemas IoT modernos. En un esquema AEAD, el algoritmo recibe como entrada un mensaje en texto plano, una clave secreta, un nonce y, opcionalmente, datos asociados que no serán cifrados pero sí autenticados. Como salida, produce un texto cifrado y una etiqueta de autenticación que permite verificar la integridad del mensaje \cite{Katz2020IntroductionCryptography}.

Esta estructura resulta adecuada para protocolos de comunicación en IoT porque permite proteger tanto la carga útil del mensaje como ciertos metadatos necesarios para el funcionamiento del protocolo. Por ejemplo, en una red de sensores puede ser necesario autenticar identificadores, contadores, encabezados o información de enrutamiento, aunque dichos campos deban permanecer visibles para los nodos intermedios. AEAD permite lograr este equilibrio al separar los datos cifrados de los datos únicamente autenticados. Si la etiqueta de autenticación no es válida, el algoritmo debe rechazar el mensaje, propiedad fundamental para evitar ataques de modificación, repetición o inyección de mensajes maliciosos. La formalización de las operaciones de cifrado y descifrado de un esquema AEAD se presenta en la Sección~\ref{sec:aead} de Marco teórico.

\subsection{ASCON como estándar de referencia}

ASCON es una familia de algoritmos criptográficos basada en una construcción tipo esponja. Fue diseñada para proporcionar cifrado autenticado y funciones hash ligeras con un balance positivo entre seguridad, rendimiento y bajo consumo de recursos \cite{Dobraunig2021AsconHashing,Bertoni2007SpongeFunctions,Dobraunig2021AsconProcess}. Su diseño se apoya en una permutación interna que opera sobre un estado de tamaño fijo y que alterna fases de absorción, transformación y extracción de salida.

La construcción tipo esponja puede entenderse como un mecanismo que recibe información de entrada, la incorpora gradualmente dentro de un estado interno y posteriormente produce una salida a partir de dicho estado. Esta aproximación resulta flexible y permite construir funciones hash, generadores pseudoaleatorios y esquemas de cifrado autenticado utilizando una misma primitiva de permutación. El detalle formal de las fases de absorción y extracción, y de la relación entre tasa y capacidad del estado interno, se presenta en la Sección~\ref{sec:ascon} de Marco teórico.

En el contexto de IoT, ASCON presenta varias ventajas. En primer lugar, proporciona cifrado y autenticación dentro de una misma construcción, reduciendo la necesidad de combinar múltiples algoritmos. En segundo lugar, su diseño es eficiente en plataformas con recursos restringidos, tanto en software como en hardware. En tercer lugar, ha sido sometido a un amplio proceso de análisis por parte de la comunidad criptográfica, lo cual fortalece la confianza en su seguridad.

No obstante, ASCON no resuelve por sí solo todos los problemas de seguridad de un sistema IoT. Como cualquier algoritmo criptográfico simétrico, su seguridad depende de la correcta generación y administración de claves, nonces y semillas. Un algoritmo AEAD robusto puede verse comprometido si el nonce se reutiliza de manera indebida, si las claves se generan con baja entropía o si el generador de números aleatorios produce secuencias predecibles. Esta observación conecta directamente la criptografía ligera con la necesidad de fuentes de entropía confiables.

\subsection{Comparación entre AES-128, SPECK y ASCON}

Diversos trabajos experimentales han comparado el comportamiento de algoritmos como AES-128, SPECK y ASCON en plataformas IoT restringidas. Estos tres algoritmos resultan particularmente útiles para el análisis porque representan enfoques de diseño diferentes.

AES-128 es un estándar ampliamente adoptado y estudiado \cite{Radhakrishnan2024EfficiencyDevices}, basado en una red de sustitución-permutación. Su principal ventaja es la madurez criptográfica y la amplia disponibilidad de implementaciones optimizadas. Sin embargo, no fue diseñado originalmente para dispositivos de bajo consumo, por lo que puede requerir optimizaciones específicas para operar eficientemente en microcontroladores.

SPECK fue diseñado con un enfoque explícitamente ligero. Su estructura ARX permite implementaciones compactas y eficientes en software, lo cual lo hace atractivo para dispositivos con restricciones severas. Sin embargo, parte de la literatura ha señalado que, aunque su desempeño es favorable, su adopción debe evaluarse cuidadosamente considerando el nivel de confianza criptográfica requerido y la sensibilidad de la aplicación.

Por su parte, ASCON representa una aproximación moderna basada en cifrado autenticado y construcción esponja. Su principal ventaja es que ofrece confidencialidad e integridad con bajo overhead, siendo especialmente adecuado para aplicaciones IoT donde se requiere proteger mensajes cortos, frecuentes y generados por dispositivos con energía limitada.

La Tabla~\ref{tab:comparacion_lwc} resume las principales características de estos algoritmos.

\begin{table}[htbp]
\centering
\caption{Comparación conceptual entre AES-128, SPECK y ASCON}
\label{tab:comparacion_lwc}
\begin{tabularx}{\textwidth}{@{}p{0.18\textwidth}p{0.25\textwidth}p{0.25\textwidth}X@{}}
\toprule
\textbf{Algoritmo} & \textbf{Estructura} & \textbf{Ventajas} & \textbf{Limitaciones} \\
\midrule
AES-128 & Red de sustitución-permutación & Alto nivel de análisis criptográfico, amplia adopción y disponibilidad de implementaciones & Puede ser costoso en dispositivos extremadamente restringidos si no cuenta con aceleración u optimización \\
SPECK & Red ARX & Bajo tamaño de código, operaciones simples y alto rendimiento en software & Menor nivel de confianza institucional frente a estándares ampliamente adoptados \\
ASCON & Construcción esponja AEAD & Proporciona confidencialidad, integridad y autenticación con bajo overhead & Depende de una correcta generación de claves y nonces para preservar sus garantías de seguridad \\
\bottomrule
\end{tabularx}
\end{table}

La comparación anterior evidencia que la selección de un algoritmo criptográfico para IoT no puede realizarse únicamente con base en su velocidad. Es necesario considerar el modelo de amenaza, la criticidad de la aplicación, la disponibilidad de memoria, el consumo energético, la facilidad de implementación segura y la calidad de la fuente de entropía utilizada para inicializar el sistema.

\subsection{Limitaciones de criptografía ligera}

Aunque la criptografía ligera constituye una respuesta adecuada a las restricciones de los dispositivos IoT, no elimina todos los riesgos presentes en estos entornos. Un algoritmo ligero puede ofrecer excelentes propiedades matemáticas y, aun así, ser vulnerable si se implementa de manera incorrecta, si no se protege contra ataques de canal lateral o si depende de claves generadas con una baja entropía \cite{MITRE2024CWE-327:Algorithm}.

En este sentido, la criptografía ligera debe entenderse como un componente dentro de una arquitectura de seguridad más amplia. Dicha arquitectura debe incluir mecanismos de arranque seguro, actualización remota protegida, autenticación de dispositivos, control de acceso, gestión segura de claves y monitoreo del comportamiento del sistema. La seguridad de un dispositivo IoT no depende exclusivamente del algoritmo de cifrado, sino de la interacción correcta entre todos estos componentes.

Por esta razón, la presente investigación no se limita a estudiar un algoritmo de cifrado ligero, sino que aborda uno de los elementos fundamentales de cualquier arquitectura criptográfica: la generación de números aleatorios. En particular, se analiza la viabilidad de utilizar QPP-RNG como fuente de entropía para fortalecer la generación de material criptográfico en dispositivos IoT críticos.

%%=========== 
%% SECCION 4
%%=========== 

\section{Generadores de números aleatorios y fuentes de entropía}

La generación de números aleatorios constituye uno de los pilares fundamentales de la criptografía moderna. Aunque frecuentemente recibe menos atención que los algoritmos de cifrado o los protocolos de autenticación, la aleatoriedad es indispensable para la generación de claves, nonces, vectores de inicialización, semillas criptográficas, desafíos de autenticación y parámetros efímeros utilizados en múltiples protocolos de seguridad.

Un sistema criptográfico puede ser matemáticamente robusto y, aun así, resultar vulnerable si utiliza números aleatorios predecibles. Esta situación ha sido documentada en diversos incidentes históricos, donde fallas en generadores de números aleatorios permitieron recuperar claves privadas, romper protocolos de autenticación o reducir significativamente el espacio efectivo de búsqueda de un atacante. Por esta razón, la calidad de la aleatoriedad debe evaluarse como una propiedad crítica de todo sistema criptográfico.

\subsection{Aleatoriedad, impredecibilidad y entropía}

En criptografía, una secuencia aleatoria no solamente debe parecer estadísticamente uniforme, sino que debe ser impredecible para un adversario. Esta diferencia es de vital importancia. Una secuencia puede superar algunas pruebas estadísticas básicas y, sin embargo, ser generada por un proceso determinístico cuya semilla sea recuperable, en cuyo caso no sería segura desde el punto de vista criptográfico \cite{Shannon1949CommunicationSystems}.

La literatura distingue principalmente dos medidas de incertidumbre: la entropía de Shannon, la cual captura la incertidumbre promedio de una fuente, y la min-entropía, que estima el peor caso frente a un adversario que intenta predecir el valor más probable \cite{Turan2018RecommendationGeneration}. Esta segunda medida resulta especialmente relevante en criptografía, pues una fuente puede presentar una entropía promedio aceptable y aun así mostrar valores individuales altamente predecibles. La formalización matemática de ambas medidas, junto con ejemplos numéricos aplicados a bytes de 8 bits, se presenta en la Sección~\ref{sec:aleatoriedad_entropia} del capítulo de Marco teórico.

\subsection{Tipos de generadores de números aleatorios}

Los generadores de números aleatorios utilizados en criptografía pueden clasificarse en tres grupos principales: generadores físicos o verdaderos (TRNG), generadores pseudoaleatorios (PRNG) y generadores determinísticos criptográficamente seguros (CSPRNG o DRBG).

Los TRNG obtienen aleatoriedad a partir de fenómenos físicos no determinísticos, como ruido térmico, osciladores de anillo o fluctuaciones eléctricas. Su principal ventaja es producir entropía real, aunque suelen requerir hardware especializado \cite{Barker2015RecommendationGenerators}. Los PRNG producen secuencias aparentemente aleatorias a partir de una semilla determinística, por lo que no son necesariamente adecuados para criptografía si dicha semilla puede recuperarse. Los CSPRNG, en cambio, están diseñados específicamente para resistir predicción adversarial, aunque su seguridad continúa dependiendo críticamente de la calidad de la semilla inicial \cite{Barker2015RecommendationGenerators}. La Tabla~\ref{tab:comparacion_rng} del capítulo de Marco teórico resume estas categorías junto con los RNG basados en \textit{jitter} de sistema, familia a la que pertenece QPP-RNG.

\subsection{Fuentes de entropía en dispositivos IoT}

En dispositivos IoT, la obtención de entropía presenta dificultades particulares. Muchos microcontroladores de bajo costo no incorporan hardware para TRNG, o lo hacen mediante componentes cuya calidad puede variar entre fabricantes \cite{Seyhan2022ClassificationTaxonomy}. Además, algunos dispositivos operan en entornos altamente determinísticos, ejecutando firmware simple, tareas periódicas y patrones de comunicación predecibles, lo cual reduce la variabilidad disponible para recolectar entropía.

A diferencia de los computadores convencionales, donde el sistema operativo recolecta entropía a partir de múltiples eventos como movimiento del mouse, interrupciones o latencias de red, un dispositivo IoT puede carecer de muchos de estos eventos. Por ejemplo, un sensor industrial puede ejecutar siempre el mismo ciclo de medición y transmisión, sin interacción humana y con una pila de software mínima \cite{Klimushyn2022Crypto-ResistantDevices}.

Esta condición obliga a estudiar fuentes alternativas de entropía, entre ellas el \textit{jitter} de temporizadores, variaciones en tiempos de ejecución y diferencias microarquitecturales. Estas fuentes no requieren necesariamente hardware especializado, pero deben evaluarse cuidadosamente para demostrar que contienen suficiente incertidumbre y que no pueden ser modeladas fácilmente por un atacante \cite{Kuang2025IID-BasedJitter}.

\subsection{Evaluación estadística de la aleatoriedad}

La evaluación de generadores de números aleatorios se apoya en baterías de pruebas estadísticas y estimadores de entropía. Entre los instrumentos más utilizados se encuentran NIST SP 800-22 \cite{Rukhin2010AApplications}, orientado a detectar patrones estadísticos en secuencias binarias; NIST SP 800-90B \cite{Turan2018RecommendationGeneration}, que en lugar de verificar apariencia de aleatoriedad busca estimar de manera conservadora la min-entropía disponible en una fuente no determinística y ENT \cite{Walker2008ENT:Program}, una herramienta más simple orientada a identificar sesgos a nivel de byte.

Superar NIST SP 800-22 no garantiza por sí solo que una fuente sea criptográficamente segura, pues estas pruebas analizan propiedades observables de la salida sin estimar cuánta entropía real contiene la fuente. Por ello, la literatura recomienda combinar ambos estándares con herramientas complementarias como ENT. El detalle metodológico de cada prueba, incluyendo la interpretación de valores-$p$ y la distinción entre fuentes IID y no-IID, se desarrolla en la Sección~\ref{sec:evaluacion_estadistica} de Marco teórico.

\subsection{Relación entre RNG y criptografía ligera}

La relación entre generadores de números aleatorios y criptografía ligera es directa. En un sistema IoT protegido mediante ASCON, AES-128 u otro algoritmo simétrico, la seguridad dependerá de la calidad de la clave, del nonce y de cualquier semilla utilizada por los mecanismos internos. Si estos valores son predecibles, repetidos o generados con baja entropía, el algoritmo puede perder sus garantías de seguridad \cite{MITRE2024CWE-327:Algorithm}.

Este riesgo es especialmente importante en esquemas AEAD, donde la reutilización de nonces bajo una misma clave puede comprometer la confidencialidad e integridad de los mensajes. Por ello, un dispositivo IoT no solamente necesita un algoritmo ligero eficiente, sino también un mecanismo confiable para generar valores únicos, impredecibles y adecuados para el contexto criptográfico.

La literatura sobre IoT ha tendido a concentrarse en la comparación de algoritmos de cifrado, midiendo variables como tiempo de ejecución, uso de RAM, ocupación de ROM, throughput y consumo energético \cite{Thakor2021LightweightOpportunities,Radhakrishnan2024EfficiencyDevices}. Aunque estas métricas son esenciales, no capturan completamente la seguridad del sistema si no se evalúa también la fuente de entropía. Esta brecha justifica el enfoque de la presente investigación, orientado a analizar QPP-RNG como componente complementario dentro de una arquitectura de criptografía ligera.

%%=========== 
%% SECCION 5
%%=========== 

\section{\textit{Quantum Permutation Pad} y QPP-RNG}

El Quantum Permutation Pad (QPP) es una propuesta criptográfica basada en el uso de espacios de permutaciones derivados de sistemas de múltiples qubits \cite{Kuang2022QuantumCryptography,Kuang2021QuantumPad}. Su fundamento conceptual consiste en extender la idea del \textit{One-Time Pad} clásico hacia un espacio de permutaciones significativamente más amplio, aprovechando propiedades como la biyectividad, reversibilidad, composición de permutaciones y no conmutatividad.

En un sistema clásico de $n$ bits, el espacio de posibles valores contiene $2^n$ estados. En cambio, el espacio de permutaciones sobre esos estados contiene $(2^n)!$ permutaciones posibles. Este crecimiento factorial genera un espacio combinatorio extremadamente grande, lo cual resulta atractivo desde el punto de vista criptográfico \cite{Kuang2020ShannonSpace}. Por ejemplo, para $n=8$, existen $2^8=256$ estados, pero el número de permutaciones posibles es $256!$, una cantidad enormemente superior al espacio clásico de claves de 8 bits.

\subsection{Fundamentos de QPP}

Una compuerta de permutación cuántica puede entenderse como una transformación biyectiva sobre la base computacional de un sistema de $n$ qubits, que transforma cada estado de entrada en un estado de salida único mediante una relación biyectiva, garantizando que la operación sea reversible e indispensable para el cifrado y descifrado \cite{Kuang2022QuantumCryptography}. En una representación matricial, cada permutación puede describirse mediante una matriz cuadrada con exactamente un valor distinto de cero en cada fila y en cada columna, lo cual permite tratarla como una operación lineal reversible sobre la base computacional.

Una de las propiedades más relevantes de QPP es la no conmutatividad: en general, dos permutaciones $\hat{p}_i$ y $\hat{p}_j$ no satisfacen $\hat{p}_i \hat{p}_j = \hat{p}_j \hat{p}_i$, de manera que el orden de aplicación afecta el resultado final \cite{Kuang2024QPPGroup}. Desde la perspectiva criptográfica, esta propiedad introduce una barrera adicional para el adversario, pues dificulta invertir la transformación sin conocer la secuencia exacta de permutaciones utilizada. La formalización matricial completa de esta propiedad, incluyendo la definición del conmutador $[P_i,P_j]$, se presenta en la Sección~\ref{sec:qpp} de Marco teórico.

\subsection{QPP como extensión del \textit{One-Time Pad}}

El \textit{One-Time Pad} clásico ofrece secreto perfecto bajo tres condiciones: la clave debe ser verdaderamente aleatoria, tener al menos la misma longitud que el mensaje y utilizarse una sola vez. Aunque estas condiciones proporcionan una seguridad teórica ideal, su cumplimiento práctico resulta difícil, especialmente por la necesidad de distribuir y almacenar claves de gran tamaño.

QPP puede interpretarse como una generalización del OTP hacia espacios de permutaciones \cite{Kuang2021QuantumPad,Kuang2022QuantumCryptography}. En lugar de operar mediante una simple suma módulo 2 entre mensaje y clave, QPP aplica transformaciones biyectivas seleccionadas dentro de un espacio combinatorio mucho más grande. Esta generalización permite explorar esquemas en los cuales la seguridad no depende únicamente del tamaño lineal de una clave binaria, sino de la estructura y selección de permutaciones dentro de un espacio factorial. El crecimiento de dicho espacio, para distintos valores de $n$, se ilustra en la Tabla~\ref{tab:crecimiento_qpp} de Marco teórico, y explica por qué QPP ha sido propuesto como una alternativa relevante para criptografía cuánticamente segura y para mecanismos de fortalecimiento de algoritmos simétricos \cite{Kuang2025QuantumCryptography}.

\subsection{QPP en criptografía ligera}

Una de las aplicaciones propuestas de QPP consiste en integrarlo con algoritmos simétricos existentes. En particular, la literatura ha explorado variantes como AES-QPP, donde ciertas operaciones tradicionales de AES son reemplazadas o complementadas mediante permutaciones seleccionadas secretamente \cite{Kuang2021QuantumPad}.

En AES tradicional, la etapa \textit{SubBytes} aplica una caja de sustitución fija, mientras que \textit{AddRoundKey} introduce entropía mediante una operación XOR con la clave de ronda. En AES-QPP, estas operaciones pueden reinterpretarse mediante matrices de permutación seleccionadas a partir de material secreto. Esta aproximación busca incrementar la confusión y reducir la linealidad entre rondas, fortaleciendo la resistencia frente a análisis diferencial y lineal \cite{Kuang2022QuantumSystems,Perepechaenko2023QuantumComputers}.

Aunque esta línea de investigación no reemplaza directamente a estándares consolidados como ASCON, resulta relevante porque muestra que QPP puede ser implementado sobre sistemas clásicos y puede interactuar con primitivas criptográficas existentes \cite{Lou2022BenchmarkPad}. Para la presente tesis, esta observación es importante porque permite considerar QPP no solamente como una propuesta teórica asociada a computación cuántica, sino como una familia de técnicas potencialmente aplicables en software convencional y dispositivos embebidos.

\subsection{QPP-RNG}

QPP-RNG, o \textit{Quantum Permutation Pad Random Number Generator}, es una propuesta orientada a la generación de números aleatorios a partir de fuentes de entropía basadas en software\cite{Kuang2025QPP-RNG:Randomness}. Su objetivo consiste en aprovechar variaciones internas del sistema computacional, como jitter temporal, comportamiento microarquitectural, acceso a memoria, tiempos de ejecución y eventos del sistema operativo, para producir secuencias con alta calidad estadística \cite{Kuang2025IID-BasedJitter}.

A diferencia de un TRNG físico, QPP-RNG no requiere necesariamente hardware especializado. Esta característica lo hace atractivo para entornos IoT, donde no todos los dispositivos incorporan fuentes físicas confiables de aleatoriedad. Al mismo tiempo, al tratarse de un mecanismo de software, su seguridad debe evaluarse rigurosamente mediante pruebas estadísticas y estimadores de entropía \cite{Vrana2025RawEvaluation}.

La ventaja conceptual de QPP-RNG radica en combinar fuentes de variabilidad del sistema con mecanismos de transformación inspirados en QPP. De esta manera, la salida no depende únicamente de una fuente simple de jitter, sino de un proceso de recolección, mezcla y transformación orientado a maximizar la impredecibilidad observable.

\subsection{Pertinencia de QPP-RNG para IoT crítico}

El uso de QPP-RNG en dispositivos IoT críticos resulta relevante por tres razones principales.

En primer lugar, muchos dispositivos IoT carecen de generadores físicos de números aleatorios o incorporan TRNG cuya calidad no siempre puede verificarse fácilmente. Una alternativa de software portable puede facilitar la generación de entropía en plataformas heterogéneas.

En segundo lugar, los sistemas críticos requieren no solamente eficiencia, sino también robustez frente a escenarios adversos. En aplicaciones como salud, transporte, energía o industria, una falla criptográfica puede generar consecuencias que van más allá de la pérdida de información, afectando procesos físicos o servicios esenciales.

En tercer lugar, la transición hacia criptografía post-cuántica y arquitecturas cuánticamente conscientes exige fortalecer todos los componentes criptográficos del sistema. Aunque los algoritmos simétricos son menos vulnerables que RSA o ECC frente a ataques cuánticos, continúan dependiendo de claves y semillas generadas con alta entropía \cite{Grover1996ASearch,Turan2018RecommendationGeneration}.

Por lo anterior, QPP-RNG constituye un objeto de estudio adecuado para esta investigación. Su evaluación permitirá determinar si puede integrarse dentro de una librería criptográfica ligera implementada en Rust y si su desempeño resulta viable frente a las restricciones propias de dispositivos IoT críticos.

%%=========== 
%% SECCION 6
%%=========== 

\section{Análisis crítico de la literatura y brecha de investigación}

La revisión del estado del arte permite identificar avances significativos en criptografía ligera, seguridad para IoT y propuestas criptográficas relacionadas con QPP. Sin embargo, también evidencia una brecha importante entre la madurez de los algoritmos de cifrado ligero y el estudio de fuentes de entropía adecuadas para dispositivos IoT críticos.

Por un lado, la literatura sobre criptografía ligera ha alcanzado un nivel considerable de consolidación. Existen múltiples familias de algoritmos diseñados para plataformas restringidas, y el proceso de estandarización liderado por NIST ha permitido identificar propuestas robustas como ASCON \cite{NationalInstituteofStandards2023AnnouncingSelection,McKay2017ReportCryptography}. De igual manera, estudios experimentales han comparado algoritmos como AES-128, SPECK y ASCON utilizando métricas de ejecución, memoria, throughput y latencia \cite{Radhakrishnan2024EfficiencyDevices,Thakor2021LightweightOpportunities}.

No obstante, gran parte de estos trabajos se concentra en el desempeño del algoritmo de cifrado, dejando en segundo plano la evaluación de la fuente de entropía utilizada para generar claves, nonces y semillas. Esta omisión resulta crítica, pues la seguridad efectiva de un sistema criptográfico depende tanto del algoritmo como del material aleatorio utilizado durante su inicialización y operación.

Por otro lado, la literatura sobre QPP muestra un marco teórico prometedor basado en espacios de permutaciones, no conmutatividad y crecimiento factorial del espacio de transformaciones. Estas propiedades han sido exploradas principalmente en el contexto de criptografía cuánticamente segura, variantes de AES-QPP y esquemas de cifrado simétrico o asimétrico \cite{Kuang2022QuantumCryptography,Kuang2024QPPGroup,Kuang2022QuantumSystems}. Sin embargo, todavía existe un número limitado de investigaciones orientadas a su implementación práctica en dispositivos IoT con restricciones severas.

En particular, se identifican cuatro brechas principales. La primera es la falta de implementaciones abiertas y modulares de QPP-RNG, particularmente en lo que concierne dispositivos IoT. Segundo, la ausencia de evaluaciones sistemáticas de QPP-RNG en términos de memoria, tiempo de ejecución, portabilidad y calidad de entropía. Tercero, la limitada integración entre QPP-RNG y algoritmos de criptografía ligera como ASCON. Por último, la escasa exploración del uso de Rust como lenguaje seguro para implementar primitivas criptográficas ligeras y fuentes de entropía en sistemas embebidos.

Estas brechas justifican el desarrollo de la presente investigación. En lugar de proponer un nuevo algoritmo de cifrado, este trabajo se enfoca en un componente transversal: la generación de entropía. La hipótesis central es que una adaptación optimizada de QPP-RNG puede constituir una fuente viable de aleatoriedad para dispositivos IoT críticos, siempre que su calidad estadística, consumo de recursos y facilidad de integración sean evaluados rigurosamente.

\subsection{Síntesis del estado del arte}

La Tabla~\ref{tab:sintesis_estado_arte} resume los principales hallazgos de la revisión bibliográfica y su relación con la presente tesis.

\begin{table}[htbp]
\centering
\caption{Síntesis del estado del arte y relación con la investigación}
\label{tab:sintesis_estado_arte}
\begin{tabularx}{\textwidth}{@{}p{0.22\textwidth}p{0.33\textwidth}X@{}}
\toprule
\textbf{Área} & \textbf{Hallazgo principal} & \textbf{Relación con la tesis} \\
\midrule
IoT crítico & Los dispositivos presentan restricciones severas de memoria, energía y procesamiento & La solución propuesta debe ser ligera y portable \\
Criptografía ligera & ASCON ofrece una referencia robusta para cifrado autenticado en dispositivos restringidos & Se utiliza como punto de comparación e integración \\
RNG criptográficos & La seguridad depende de fuentes de entropía confiables, no solamente del algoritmo de cifrado & QPP-RNG se estudia como fuente alternativa de entropía \\
QPP & El espacio de permutaciones ofrece propiedades criptográficas atractivas, incluyendo no conmutatividad y crecimiento factorial & Proporciona el fundamento conceptual de QPP-RNG \\
Rust & Su modelo de memoria favorece implementaciones seguras de software crítico & Se selecciona como lenguaje para la librería propuesta \\
\bottomrule
\end{tabularx}
\end{table}

En conclusión, el estado del arte muestra que existe una necesidad concreta de estudiar fuentes de entropía compatibles con criptografía ligera e IoT crítico. Aunque los avances en algoritmos como ASCON han fortalecido la protección de dispositivos restringidos, la generación segura de números aleatorios continúa siendo un componente esencial y menos explorado. Esta tesis aborda dicha necesidad mediante el diseño, implementación y evaluación de una librería en Rust que incorpora QPP-RNG como fuente de entropía para criptografía ligera.
